% !TEX program = xelatex
% NOTE: Для корректного Unicode используйте XeLaTeX (или LuaLaTeX).
\documentclass[11pt]{article}
\usepackage{fontspec}
\setmainfont{CMU Serif}
\usepackage[margin=0.75in]{geometry}
\usepackage{amsmath}
\usepackage{enumitem}
\usepackage{hyperref}

\setlength{\parindent}{0pt}
\setlength{\parskip}{0.32\baselineskip}
\setlist[itemize]{leftmargin=1.2em, topsep=2pt, itemsep=1pt, parsep=0pt, partopsep=0pt}
\setlist[enumerate]{leftmargin=1.45em, topsep=2pt, itemsep=1pt, parsep=0pt, partopsep=0pt}
\setlength{\abovedisplayskip}{5pt}
\setlength{\belowdisplayskip}{5pt}
\setlength{\abovedisplayshortskip}{3pt}
\setlength{\belowdisplayshortskip}{3pt}

\begin{document}

\begin{center}
{\LARGE \textbf{ЕГЭ Информатика --- Задание 13}}\\[-2pt]
{\large \textbf{Сети и IP-адресация}}
\end{center}

\section*{Теория}

Задание №13 проверяет понимание \textbf{основ сетей и IP-адресации}. Обычно это базовое задание, на которое уходит около \textbf{$\approx$ 3 минут}, а при хорошем знании теории --- \textbf{10--30 секунд}.

\textbf{Что нужно знать:}
\begin{enumerate}
  \item \textbf{IP-адрес}
  \item \textbf{маска сети}
  \item \textbf{сеть и её компоненты}
\end{enumerate}

\subsection*{IP-адрес}

\textbf{IP-адрес} --- уникальный идентификатор устройства в сети. Он состоит из 4 байтов, то есть из 32 бит. Каждый байт принимает значение от 0 до 255.

\[
1 \text{ байт} = 8 \text{ бит}, \qquad 4 \text{ байта} = 32 \text{ бита}
\]

Пример IP-адреса: \texttt{192.168.42.17}. В двоичной записи этот адрес выглядит так: \texttt{11000000.10101000.00101010.00010001}.

\textbf{Алгоритм перевода в двоичную систему:}
\begin{enumerate}
  \item каждый байт переводится отдельно
  \item если битов меньше 8, добавляются ведущие нули
  \item байты записываются через точку в исходном порядке
\end{enumerate}

\subsection*{Маска сети}

\textbf{Маска сети} определяет, какая часть IP-адреса относится к сети, а какая --- к узлу. Она имеет ту же структуру: 4 байта, то есть 32 бита.

Главное свойство маски: в двоичной записи сначала идут единицы, потом нули. Например,
\[
11111111111111100000000000000000
\]
корректно, а
\[
111100010000
\]
некорректно, потому что после нуля снова появляется единица.

\textbf{Примеры корректных масок:}
\begin{itemize}
  \item \texttt{255.255.255.0}
  \item \texttt{255.255.224.0}
  \item \texttt{255.255.255.255}
  \item \texttt{0.0.0.0}
\end{itemize}

Некорректный пример: \texttt{11111111.11111111.00111101.00000000}.

\subsection*{CIDR-запись}

Маску можно записывать кратко: \texttt{/n}, где $n$ --- количество единиц в маске.

Например, маска \texttt{255.255.255.0} в двоичном виде равна \texttt{11111111.11111111.11111111.00000000}, значит, в ней 24 единицы, то есть это \texttt{/24}.

\subsection*{Компоненты сети}

Любая сеть содержит:
\begin{enumerate}
  \item маску
  \item адрес сети
  \item широковещательный адрес
  \item узлы
\end{enumerate}

\textbf{Адрес сети} --- первый IP-адрес сети. \textbf{Broadcast} --- последний IP-адрес сети. \textbf{Узлы} --- все остальные устройства сети.

\[
\text{адрес сети} \neq \text{узел}, \qquad \text{broadcast} \neq \text{узел}
\]

Поэтому
\[
\text{количество узлов} = \text{количество IP} - 2
\]

\subsection*{Количество адресов и узлов}

Если $n$ --- количество единиц в маске, то число нулей:
\[
k = 32 - n
\]

Тогда:
\[
\text{количество IP-адресов} = 2^k, \qquad \text{количество узлов} = 2^k - 2
\]

Пример: маска \texttt{255.255.240.0} имеет вид \texttt{11111111.11111111.11110000.00000000}. В ней 20 единиц, значит,
\[
k = 32 - 20 = 12, \qquad 2^{12} = 4096, \qquad 4096 - 2 = 4094
\]

\textbf{Главное правило сети:}
\[
\text{первый адрес} = \text{адрес сети}, \qquad \text{последний адрес} = \text{broadcast}
\]

Все остальные адреса --- узлы.

\subsection*{Python для задания 13}

Ручной перевод IP и побитовые операции занимают слишком много времени, поэтому обычно используют модуль \texttt{ipaddress}.

\begin{verbatim}
from ipaddress import *
net = ip_network("IP/MASK", False)
\end{verbatim}

Полезные свойства:
\begin{verbatim}
net.network_address
net.broadcast_address
net.netmask
\end{verbatim}

Если маска неизвестна, можно перебирать все варианты:
\begin{verbatim}
for m in range(33):
    net = ip_network(f"148.195.140.28/{m}", False)
    print(net)
\end{verbatim}

\section*{Решение прототипов}

\subsection*{1. Найти адрес сети}

\textbf{Условие.} По заданным IP-адресу узла сети и маске определите адрес сети: IP-адрес: \texttt{10.8.248.131}, маска: \texttt{255.255.224.0}.

\textbf{Решение.} Адрес сети получается побитовым И между IP-адресом и маской. Удобнее всего решать через Python:
\begin{verbatim}
from ipaddress import *

net = ip_network(f"10.8.248.131/255.255.224.0", False)
print(net)
print(net[0])
print(net.network_address)
\end{verbatim}

\textbf{Что делает этот код:}
\begin{itemize}
  \item \texttt{ip\_network(..., False)} создаёт объект сети по IP-адресу и маске
  \item параметр \texttt{False} разрешает передавать не адрес сети, а адрес любого узла из этой сети
  \item \texttt{print(net)} печатает всю сеть в сокращённом виде: \texttt{10.8.224.0/19}
  \item \texttt{print(net[0])} печатает первый адрес сети, а первый адрес всегда является адресом сети
  \item \texttt{print(net.network\_address)} обращается к специальному свойству объекта сети и тоже даёт адрес сети
\end{itemize}

То есть две последние строки выводят одно и то же значение. В этой задаче:
\[
net[0] = net.network\_address = 10.8.224.0
\]

Получаем \texttt{10.8.224.0}.

\textbf{Ответ:} \texttt{10.8.224.0}.

\subsection*{2. Найти байт маски по IP и адресу сети}

\textbf{Условие.} Для узла с IP-адресом \texttt{220.128.112.142} адрес сети равен \texttt{220.128.96.0}. Чему равен третий слева байт маски? Ответ запишите в виде десятичного числа.

\textbf{Решение.} Эту задачу удобно решать перебором всех возможных масок и смотреть, при какой маске получится нужный адрес сети.

\begin{verbatim}
from ipaddress import *

for m in range(33):
    net = ip_network(f"220.128.112.142/{m}", False)
    print(net, net.netmask)
\end{verbatim}

\textbf{Как это работает:}
\begin{itemize}
  \item \texttt{m} --- это количество единиц в маске, то есть префикс от \texttt{/0} до \texttt{/32}
  \item для каждого значения строится сеть, в которую попадает адрес \texttt{220.128.112.142}
  \item \texttt{print(net, net.netmask)} печатает сразу адрес сети и полную маску
  \item дальше нужно просто найти глазами строку, где адрес сети равен \texttt{220.128.96.0}
\end{itemize}

В выводе появится строка вида:
\begin{verbatim}
220.128.96.0/19 255.255.224.0
\end{verbatim}

Значит, нужная маска --- \texttt{255.255.224.0}. Теперь берём третий слева байт этой маски:
\[
255.255.224.0 \Rightarrow 224
\]

То есть здесь даже не обязательно вручную переводить байты в двоичную систему: достаточно перебрать маски, найти нужный адрес сети и считать ответ из найденной маски.

\textbf{Ответ:} \texttt{224}.

\subsection*{3. Найти минимальное количество единиц в маске}

\textbf{Условие.} Для узла с IP-адресом \texttt{148.195.140.28} адрес сети равен \texttt{148.195.140.0}. Найдите наименьшее возможное количество единиц в двоичной записи маски подсети.

\textbf{Решение.} Эту задачу тоже удобно решать перебором всех возможных масок.

\begin{verbatim}
from ipaddress import *

for m in range(33):
    net = ip_network(f"148.195.140.28/{m}", False)
    print(net)
\end{verbatim}

\textbf{Как рассуждать по выводу:}
\begin{itemize}
  \item программа перебирает все префиксы от \texttt{/0} до \texttt{/32}
  \item для каждого префикса она печатает адрес сети
  \item нужно найти все строки, где сеть равна \texttt{148.195.140.0}
  \item из них выбрать \textbf{минимальное} значение \texttt{m}, потому что в задаче спрашивают минимальное количество единиц в маске
\end{itemize}

В выводе впервые встретится строка:
\begin{verbatim}
148.195.140.0/22
\end{verbatim}

Это значит, что маска \texttt{/22} уже подходит. Все более длинные маски тоже могут давать тот же адрес сети, но нам нужна именно минимальная.

То есть здесь не обязательно вручную переводить байты в двоичную систему: достаточно перебрать маски и взять первое подходящее значение префикса.

\textbf{Ответ:} \texttt{22}.

\subsection*{4. Два узла в одной сети: найти максимальный байт маски}

\textbf{Условие.} Два узла, находящиеся в одной сети, имеют IP-адреса \texttt{112.117.107.70} и \texttt{112.117.121.80}. Укажите наибольшее возможное значение третьего слева байта маски сети. Ответ запишите в виде десятичного числа.

\textbf{Решение.} Это удобно решать перебором всех масок. Идея такая: для каждой маски строим сеть для первого и второго IP, потом проверяем три условия:
\begin{itemize}
  \item первый IP действительно является узлом, а не адресом сети и не broadcast
  \item второй IP тоже является узлом
  \item обе сети совпадают
\end{itemize}

\begin{verbatim}
from ipaddress import *

ip1 = ip_address('112.117.107.70')
ip2 = ip_address('112.117.121.80')

for m in range(33):
    net1 = ip_network(f'{ip1}/{m}', False)
    net2 = ip_network(f'{ip2}/{m}', False)
    if net1[0] < ip1 < net1[-1]:
        if net2[0] < ip2 < net2[-1]:
            if net1 == net2:
                print(net1.netmask)
\end{verbatim}

\textbf{Что проверяет этот код:}
\begin{itemize}
  \item \texttt{net1[0] < ip1 < net1[-1]} означает, что \texttt{ip1} строго между первым и последним адресом сети
  \item значит, \texttt{ip1} точно не является ни адресом сети, ни широковещательным адресом
  \item аналогичная проверка делается для \texttt{ip2}
  \item \texttt{net1 == net2} означает, что оба адреса попали в одну и ту же сеть
  \item программа печатает все подходящие маски
\end{itemize}

Из напечатанных масок нужно взять ту, у которой третий байт максимален. Последняя подходящая маска будет:
\begin{verbatim}
255.255.224.0
\end{verbatim}

Берём третий слева байт:
\[
255.255.224.0 \Rightarrow 224
\]

То есть твоя дополнительная проверка правильная: мы заранее убеждаемся, что оба IP --- это именно узлы сети, а не её границы.

\textbf{Ответ:} \texttt{224}.

\subsection*{5. Найти количество узлов по маске}

\textbf{Условие.} Для некоторой подсети используется маска \texttt{255.255.255.128}. Сколько различных адресов компьютеров теоретически допускает эта маска, если два адреса (адрес сети и широковещательный) не используют?

\textbf{Решение.} Это можно посчитать формулой, а можно совсем быстро через Python:
\begin{verbatim}
from ipaddress import *

net = ip_network(f'0.0.0.0/255.255.255.128', False)
print(net.num_addresses - 2)
\end{verbatim}

\textbf{Почему это работает:}
\begin{itemize}
  \item здесь важна только маска, а конкретный адрес можно взять любой удобный, например \texttt{0.0.0.0}
  \item объект \texttt{net} создаётся по этой маске
  \item \texttt{net.num\_addresses} --- это общее количество адресов в такой сети
  \item из этого количества вычитаем 2, потому что нельзя использовать адрес сети и broadcast
\end{itemize}

Для маски \texttt{255.255.255.128} программа получает:
\[
128 - 2 = 126
\]

Если решать вручную, то \texttt{255.255.255.128} --- это маска \texttt{/25}, значит, остаётся 7 бит под узлы, всего адресов
\[
2^7 = 128
\]
и тогда
\[
128 - 2 = 126
\]

\textbf{Ответ:} \texttt{126}.

\subsection*{6. Сколько существует подходящих масок}

\textbf{Условие.} Для узла c IP-адресом \texttt{175.122.80.13} адрес подсети равен \texttt{175.122.80.0}. Сколько существует различных возможных значений маски, если известно, что в этой сети не менее 60 узлов? Ответ запишите в виде десятичного числа.

\textbf{Решение.} Это удобно решать перебором всех масок с дополнительной проверкой на количество узлов:
\begin{verbatim}
from ipaddress import *

for m in range(33):
    net = ip_network(f'175.122.80.13/{m}', False)
    if net.num_addresses - 2 >= 60:
        print(net)
\end{verbatim}

\textbf{Как рассуждать по этому коду:}
\begin{itemize}
  \item перебираются все возможные префиксы от \texttt{/0} до \texttt{/32}
  \item для каждого префикса строится сеть, в которую попадает адрес \texttt{175.122.80.13}
  \item условие \texttt{net.num\_addresses - 2 >= 60} оставляет только те сети, где не меньше 60 узлов
  \item дальше из вывода нужно взять только те строки, где адрес сети равен \texttt{175.122.80.0}
\end{itemize}

В выводе останутся такие подходящие сети:
\begin{verbatim}
175.122.80.0/20
175.122.80.0/21
175.122.80.0/22
175.122.80.0/23
175.122.80.0/24
175.122.80.0/25
175.122.80.0/26
\end{verbatim}

Их 7 штук, значит, существует 7 подходящих масок.

Здесь удобно именно считать количество подходящих строк в выводе.

\textbf{Ответ:} \texttt{7}.

\subsection*{7. Найти номер компьютера в сети}

\textbf{Условие.} Если маска подсети \texttt{255.255.252.0} и IP-адрес компьютера в сети \texttt{156.132.15.138}, то чему будет равен номер компьютера в сети?

\textbf{Решение.} Это можно находить через смещение от адреса сети, но можно и просто перебрать все адреса сети по порядку:
\begin{verbatim}
from ipaddress import *

net = ip_network('156.132.15.138/255.255.252.0', False)

n = 0
for x in net:
    print(n, x)
    n += 1
\end{verbatim}

\textbf{Что делает этот код:}
\begin{itemize}
  \item создаётся сеть, в которую входит IP-адрес \texttt{156.132.15.138}
  \item переменная \texttt{n} --- это номер адреса внутри сети
  \item цикл \texttt{for x in net} перебирает все адреса сети подряд: от адреса сети до broadcast
  \item команда \texttt{print(n, x)} печатает номер и сам IP-адрес
\end{itemize}

Дальше нужно просто найти в выводе строку, где будет адрес \texttt{156.132.15.138}. Для него программа напечатает:
\begin{verbatim}
906 156.132.15.138
\end{verbatim}

Значит, номер этого компьютера в сети равен \texttt{906}.

Почему это работает: перебор начинается с адреса сети, у него номер 0, затем номер увеличивается на 1 для каждого следующего адреса. Поэтому когда цикл доходит до нужного IP, число \texttt{n} и есть его номер в этой сети.

\textbf{Ответ:} \texttt{906}.

\subsection*{8. Найти наибольший IP-адрес компьютера в сети}

\textbf{Условие.} Сеть задана IP-адресом одного из входящих в неё узлов \texttt{218.194.82.148} и сетевой маской \texttt{255.255.255.192}. Найдите наибольший IP-адрес в данной сети, который может быть назначен компьютеру. В ответе укажите найденный IP-адрес без разделителей.

\textbf{Решение.} Это удобно сделать совсем коротко через объект сети:
\begin{verbatim}
from ipaddress import *

net = ip_network('218.194.82.148/255.255.255.192', False)
print(net[-2])
\end{verbatim}

\textbf{Что делает этот код:}
\begin{itemize}
  \item создаётся сеть, в которую входит адрес \texttt{218.194.82.148}
  \item объект \texttt{net} хранит все адреса сети по порядку
  \item \texttt{net[-1]} --- это последний адрес сети, то есть broadcast
  \item \texttt{net[-2]} --- это адрес перед broadcast, а значит, максимальный адрес, который можно выдать компьютеру
\end{itemize}

То есть логика здесь очень простая:
\[
\text{последний адрес} = \text{broadcast}, \qquad \text{предпоследний адрес} = \text{последний допустимый узел}
\]

Программа выводит:
\begin{verbatim}
218.194.82.190
\end{verbatim}

По условию ответ нужно записать без точек, поэтому получаем:
\[
21819482190
\]

\textbf{Ответ:} \texttt{21819482190}.

\subsection*{9. Сколько адресов удовлетворяют условию по двоичной записи}

\textbf{Условие.} Сеть задана IP-адресом \texttt{214.96.0.0} и маской сети \texttt{255.240.0.0}. Сколько в этой сети IP-адресов, у которых количество нулей в двоичной записи IP-адреса кратно трём?

\textbf{Решение.} Маска \texttt{255.240.0.0} соответствует префиксу \texttt{/12}. Значит, первые 12 бит фиксированы, остальные 20 бит перебираются.

Такие задачи удобнее всего решать программно: перебрать все адреса сети, перевести адрес в 32-битную строку и посчитать число нулей.

\begin{verbatim}
from ipaddress import *
net = ip_network("214.96.0.0/255.240.0.0")
count = 0
for ip in net:
    b = bin(int(ip))[2:].zfill(32)
    if b.count('0') % 3 == 0:
        count += 1
print(count)
\end{verbatim}

Программа даёт значение \texttt{349526}.

\textbf{Ответ:} \texttt{349526}.

\subsection*{10. Сумма единиц в двоичной записи четвёртого байта}

\textbf{Условие.} Сеть задана IP-адресом \texttt{123.222.111.192} и маской сети \texttt{255.255.255.248}. Сколько в этой сети IP-адресов, для которых сумма единиц в двоичной записи четвёртого байта IP-адреса не делится без остатка на 3? В ответе укажите только число.

\textbf{Решение.} Эту задачу удобно решить полным перебором всех адресов сети:
\begin{verbatim}
from ipaddress import *

net = ip_network(f'123.222.111.192/255.255.255.248', 0)
k = 0
for x in net:
    # ip = bin(int(x))[2:].zfill(32)
    ip = f'{int(x):032b}'
    if ip[-8:].count('1') % 3 != 0:
        k += 1
print(k)
\end{verbatim}

\textbf{Что делает этот код:}
\begin{itemize}
  \item создаётся сеть по заданному IP и маске
  \item переменная \texttt{k} считает количество подходящих адресов
  \item цикл перебирает все IP-адреса этой сети
  \item \texttt{int(x)} переводит IP-адрес в целое число
  \item \texttt{f'\{int(x):032b\}'} записывает это число как двоичную строку длиной ровно 32 бита
  \item \texttt{ip[-8:]} берёт последние 8 бит, то есть четвёртый байт IP-адреса
  \item \texttt{count('1')} считает, сколько единиц в этом байте
  \item если это количество не делится на 3, увеличиваем счётчик
\end{itemize}

Здесь сеть имеет маску \texttt{/29}, значит, в ней всего 8 адресов. Программа просто проверяет каждый из них и считает подходящие.

В результате получаем:
\begin{verbatim}
5
\end{verbatim}

\textbf{Ответ:} \texttt{5}.

\end{document}












